% SEGMENT OLP-0578-S01
% Part: set-theory
% Chapter: ord-arithmetic
% Section: multiplication

\documentclass[../../../include/open-logic-section]{subfiles}

\begin{document}

\olfileid{sth}{ord-arithmetic}{mult}
\olsection{வரிசையெண் பெருக்கல்}

இப்போது வரிசையெண் பெருக்கலை நோக்குவோம்; வரிசையெண் கூட்டலைப் போலவே
இதையும் அணுகுகிறோம். $\alpha$ வை $\beta$ ஆல் பெருக்க விரும்புகிறோம்
என்று கொள்வோம். அகலம் $\alpha$, உயரம் $\beta$ உள்ள செவ்வகக்
கட்டத்தை மனத்தில் கொள்ளுங்கள். ஒவ்வொரு கிடைவரிசையிலும் நகர்ந்து,
பின்னர் அடுத்த கிடைவரிசைக்குச் சென்று\ldots இவ்வாறு முழுக் கட்டத்தையும்
கடந்து செல்லும்போது கிடைப்பதே $\alpha$, $\beta$ ஆகியவற்றின்
பெருக்கற்பலன். வேறுவிதமாகச் சொன்னால், $\beta$ விலுள்ள
\emph{ஒவ்வொரு} உறுப்புக்கும் பதிலாக $\alpha$ வின் ஒரு பிரதியை
வைப்பதன் மூலம் $\alpha$, $\beta$ ஆகியவற்றின் பெருக்கற்பலன் கிடைக்கிறது.

இதனை முறைப்படுத்த, $\alpha$, $\beta$ ஆகியவற்றின் கார்டீசியப்
பெருக்கலின் மீது தலைகீழ் அகராதி வரிசையைப் பயன்படுத்துகிறோம்:

\begin{defn}
எந்த வரிசையெண்கள் $\alpha, \beta$ க்கும் அவற்றின் பெருக்கற்பலன்
$\alpha \ordtimes \beta = \ordtype{\alpha \times \beta, \rlexless}$.
\end{defn}
\noindent
இந்த வரையறை முறையாக அமைந்துள்ளது என்பதை மீண்டும் உறுதிசெய்ய வேண்டும்:

\begin{lem}\ollabel{ordtimeslessiswo}
எந்த வரிசையெண்கள் $\alpha$, $\beta$ க்கும்
$\tuple{\alpha \times \beta, \rlexless}$ ஒரு நன்கு
வரிசைப்படுத்தலாகும்.
\end{lem}

\begin{proof}
\olref[add]{ordsumlessiswo} இன் நிறுவலைப் போலவே.
\end{proof}
\noindent
பெருக்கல் பின்வருமாறு செயல்படுகிறது என்பதையும் எளிதில் நிறுவலாம்:

\begin{lem}\ollabel{ordtimesrecursion}
எந்த வரிசையெண்கள் $\alpha, \beta$ க்கும்:
\begin{align*}
  \alpha \ordtimes 0 &= 0\\
  \alpha \ordtimes (\beta \ordplus 1) &=
    (\alpha \ordtimes \beta) \ordplus \alpha\\
  \alpha  \ordtimes \beta &=
    \sup_{\delta < \beta}(\alpha \ordtimes \delta) &&
    \text{$\beta$ ஓர் எல்லை வரிசையெண் எனில்}.
\end{align*}
% SOURCE CORRECTION TA-STH-027: ordinary supremum handles zero left factor.
\end{lem}

\begin{proof}
பயிற்சியாக விடப்படுகிறது.
\end{proof}

% SEGMENT OLP-0578-S02
உண்மையில், கூட்டலைப் போலவே, நேரடி வரையறைக்குப் பதிலாக இந்த
மீள்வரையறைச் சமன்பாடுகளைக் கொண்டே வரிசையெண் பெருக்கலை
வரையறுத்திருக்கலாம். கூட்டலைப் போலவே, இதன் சில பண்புகளும்
பழக்கமானவையே:

\begin{lem}\ollabel{ordinalmultiplicationisnice}
$\alpha, \beta, \gamma$ வரிசையெண்கள் எனில்:
\begin{enumerate}
  \item\ollabel{ordtimes1} $\alpha \neq 0$ மற்றும்
  $\beta < \gamma$ எனில்,
  $\alpha \ordtimes \beta < \alpha \ordtimes \gamma$;
  \item\ollabel{ordtimes2} $\alpha \neq 0$ மற்றும்
  $\alpha \ordtimes \beta = \alpha\ordtimes\gamma$ எனில்,
  $\beta = \gamma$;
  \item\ollabel{ordtimes3}
  $\alpha \ordtimes (\beta \ordtimes \gamma) =
  (\alpha \ordtimes \beta) \ordtimes \gamma$;
  \item\ollabel{ordtimes4} $\alpha \leq \beta$ எனில்,
  $\alpha \ordtimes \gamma \leq \beta \ordtimes \gamma$;
  \item\ollabel{ordtimes5}
  $\alpha \ordtimes (\beta \ordplus \gamma) =
  (\alpha \ordtimes \beta )\ordplus (\alpha\ordtimes \gamma)$.
\end{enumerate}
\end{lem}

\begin{proof}
பயிற்சியாக விடப்படுகிறது.
\end{proof}

% SEGMENT OLP-0578-S03
விரும்பினால் வேறு முடிவுகளையும் நிறுவலாம் (அல்லது தேடிப் பார்க்கலாம்).
ஆனால் \olref[ord-arithmetic][add]{ordsumnotcommute} ஐக்
கருதினால், பின்வருவது ஆச்சரியமாக இருக்கக்கூடாது:

\begin{prop}
வரிசையெண் பெருக்கல் பரிமாற்றுப் பண்புடையதல்ல:
$2 \ordtimes \omega  = \omega < \omega \ordtimes 2$.
\end{prop}

\begin{proof}
$2 \ordtimes \omega = \supstrict_{n < \omega}(2\ordtimes  n)
= \omega \in \supstrict_{n < \omega}(\omega \ordplus n)
= \omega \ordplus \omega = \omega \ordtimes 2$.
\end{proof}
\noindent
இதற்கான உள்ளுணர்வுக் காரணமும் நேரடியானது. $2 \ordtimes \omega$
வைக் கணக்கிட, ஒவ்வோர் இயல் எண்ணுக்கும் பதிலாக இரண்டு பொருள்களை
வைக்கிறோம். இயல் எண்களுக்கு இடையே எல்லா ``அரை'' எண்களையும்
செருகினாலும் அதே வரிசை வகை கிடைக்கும்; அதாவது,
$\Setabs{\nicefrac{n}{2}}{n \in \omega}$ இன் இயல்பான வரிசையைக்
கருதலாம். அவ்வாறு பார்த்தால் அதன் வரிசை வகை $\omega$ வின்
வரிசை வகையே என்பது தெளிவு. ஆனால் $\omega \ordtimes 2$ வைக்
கணக்கிட, $\omega$ வின் இரண்டு பிரதிகளை ஒன்றன் பின் ஒன்றாக
வைக்கிறோம்.

\begin{prob}
\olref[sth][ord-arithmetic][mult]{ordtimeslessiswo},
\olref[sth][ord-arithmetic][mult]{ordtimesrecursion},
மற்றும்
\olref[sth][ord-arithmetic][mult]{ordinalmultiplicationisnice}
ஆகியவற்றை நிறுவுக.
\end{prob}

\end{document}
